\documentclass[aspectratio=169,11pt]{beamer}

\usetheme{Madrid}
\usecolortheme{default}
\setbeamertemplate{navigation symbols}{}

\definecolor{PKURed}{RGB}{116,17,40}
\definecolor{PKULight}{RGB}{245,236,239}

\setbeamercolor{structure}{fg=PKURed}
\setbeamercolor{title}{fg=white,bg=PKURed}
\setbeamercolor{frametitle}{fg=white,bg=PKURed}
\setbeamercolor{palette primary}{fg=white,bg=PKURed}
\setbeamercolor{palette secondary}{fg=white,bg=PKURed!88!black}
\setbeamercolor{palette tertiary}{fg=white,bg=PKURed!76!black}
\setbeamercolor{palette quaternary}{fg=white,bg=black}
\setbeamercolor{block title}{fg=white,bg=PKURed}
\setbeamercolor{block body}{fg=black,bg=PKULight}

\usepackage[utf8]{inputenc}
\usepackage[T1]{fontenc}
\usepackage{amsmath,amssymb,bm}
\usepackage{booktabs}

\title[Original-Drug Innovation and MAH]{Original-Drug Innovation, Implementation-Route Assignment,\\ and the MAH System in China}
\author{Danni Yangchen \and Wenxiao Dong \and Xinhao Wang}
\date{\today}

\begin{document}

%------------------------------------------------
\begin{frame}
  \titlepage
\end{frame}

%------------------------------------------------
\begin{frame}{1. Main Question}
\begin{itemize}
    \item This paper asks a narrow but two-layer question:
    \[
    \text{Does MAH raise original-drug innovation by changing the implementation route from ideas to products?}
    \]
    \item Mechanism object: \textbf{implementation-route assignment}.
    \item Final result object: \textbf{realized original-drug innovation}.
    \item We do not model MAH as a generic pro-innovation policy or a direct scientific productivity shock.
\end{itemize}

\vspace{0.25cm}
\begin{block}{Object hierarchy}
    \[
    \text{Mechanism: assignment} \;\Rightarrow\; \text{Result: original-drug realization}
    \]
\end{block}
\end{frame}

%------------------------------------------------
\begin{frame}{2. Motivation: Why Implementation Is the Bottleneck}
\begin{itemize}
    \item In pharmaceuticals, invention and implementation need not happen inside the same firm.
    \item A promising original-drug idea may remain unrealized if the innovator cannot retain authorization while relying on qualified external production.
    \item Before MAH, a research-originating firm often faced a difficult choice:
    \begin{enumerate}
        \item integrate downstream production,
        \item transfer effective control to an integrated producer, or
        \item leave the project unrealized.
    \end{enumerate}
    \item This means the bottleneck is often \textbf{organizational implementation}, not only scientific discovery.
\end{itemize}

\vspace{0.3cm}
\textbf{Core intuition:} MAH matters if it changes which implementation route from idea to marketable product is feasible.
\end{frame}

%------------------------------------------------
\begin{frame}{3. What MAH Changes Institutionally}
\begin{itemize}
    \item MAH decouples, in legal form, \textbf{authorization holding} from \textbf{in-house production}.
    \item A holder may retain authorization while commissioning production to a qualified outside manufacturer.
    \item But this is not simple outsourcing:
    \begin{itemize}
        \item the relationship is formalized,
        \item the holder still bears meaningful responsibility,
        \item quality and compliance obligations remain central.
    \end{itemize}
    \item So MAH does not remove regulation; it creates a more workable \textbf{external implementation route}.
\end{itemize}

\vspace{0.3cm}
\textbf{Takeaway:} MAH changes the organizational feasibility of implementation under authorization-production separation.
\end{frame}

%------------------------------------------------
\begin{frame}{4. Policy Primitive and Identification Boundary}
\begin{block}{Policy primitive}
We model MAH as lowering the route-specific friction of external implementation:
\[
\tau_{ext}(1) < \tau_{ext}(0).
\]
\end{block}

\begin{itemize}
    \item \(\tau_{ext}\) is not demand, not scientific productivity, and not a universal policy wedge.
    \item It is the governance and compliance friction of external implementation while holder responsibility is retained.
    \item Empirically, we can map outcomes to this mechanism, but we do not claim to directly observe \(\tau_{ext}\).
\end{itemize}

\vspace{0.3cm}
If MAH lowers \(\tau_{ext}\), then external implementation becomes more attractive.
\end{frame}

%------------------------------------------------
\begin{frame}{5. Baseline Model Architecture}
\begin{itemize}
    \item Time is discrete.
    \item Firms are heterogeneous with state \(z\), and choose across two project directions:
    \[
    k\in\{G,O\}\quad (\text{generic},\;\text{original-drug}).
    \]
    \item Three organizational \textbf{modes} in the baseline:
    \begin{enumerate}
        \item \textbf{Type A:} integrated mode
        \item \textbf{Type B:} research-specialized mode
        \item \textbf{Type C:} manufacturing-specialized mode
    \end{enumerate}
    \item The model includes:
    \begin{itemize}
        \item endogenous entry and exit,
        \item organizational-state switching,
        \item a contract-manufacturing market,
        \item implementation-route assignment for original-drug projects.
    \end{itemize}
\end{itemize}

\vspace{0.25cm}
\textbf{Important design choice:} MAH acts on implementation value, then feeds back into innovation direction.
\end{frame}

%------------------------------------------------
\begin{frame}{6. Economic Roles of the Three Modes}
\begin{columns}[T]
\column{0.32\textwidth}
\begin{block}{Type A}
Integrated mode
\begin{itemize}
    \item generate ideas
    \item implement internally
    \item benchmark internal route
\end{itemize}
\end{block}

\column{0.32\textwidth}
\begin{block}{Type B}
Research-specialized mode
\begin{itemize}
    \item generate ideas
    \item do not rely on full internal manufacturing for the marginal project
    \item central to the mechanism
\end{itemize}
\end{block}

\column{0.32\textwidth}
\begin{block}{Type C}
Manufacturing-specialized mode
\begin{itemize}
    \item provide qualified external production
    \item serve contract-manufacturing market
\end{itemize}
\end{block}
\end{columns}

\vspace{0.25cm}
\textbf{Interpretation:} Type B is where MAH should matter most for the marginal original-drug project.
\end{frame}

%------------------------------------------------
\begin{frame}{7. A/B/C Are Organizational States, Not Literal Firm Identities}
\begin{itemize}
    \item In reality, the same company may behave like \textbf{type B} for one drug and like \textbf{type A} for another.
    \item So in the baseline, \(A/B/C\) are \textbf{economic organizational states}, not permanent legal or accounting identities of the whole firm.
    \item The baseline abstracts from within-firm product-level heterogeneity and uses \(A/B/C\) as \textbf{dominant modes relevant for the marginal original-drug project}.
    \item This is why the paper's first-order object is \textbf{project-level implementation-route assignment}, not literal firm labels.
\end{itemize}

\vspace{0.25cm}
\begin{block}{Why keep this baseline?}
It keeps the model tractable for calibration and estimation now, while leaving open a richer multi-product extension later.
\end{block}
\end{frame}

%------------------------------------------------
\begin{frame}{8. Timing and the Core State Variable}
\begin{itemize}
    \item Each incumbent starts the period with organizational state \(s\in\{A,B,C\}\) and capability \(z\).
    \item Static production choices are optimized out; dynamic decisions are central.
    \item Within-period dynamic sequence:
    \begin{enumerate}
        \item choose project direction and effort,
        \item assign original-drug projects to implementation routes,
        \item choose stay / switch / exit.
    \end{enumerate}
    \item Baseline uses one-dimensional state:
    \[
    z\in Z\subset\mathbb{R}_+.
    \]
\end{itemize}

\vspace{0.25cm}
\textbf{Why this baseline?} Tractable enough for estimation now, expandable to richer states or product-level heterogeneity later.
\end{frame}

%------------------------------------------------
\begin{frame}{9. Direction Choice under the Type-B Organizational State}
\small
Firms operating in the type-B organizational state allocate effort across two directions:
\[
\Pi_B(z;M,p_m)
=
-f_B+
\max_{x_B^G,x_B^O\ge 0}
\left\{
x_B^G\Psi_B^G(z)
+
x_B^O\Psi_B^O(z;p_m,M)
-
C_B(x_B^G,x_B^O)
\right\}.
\]

\begin{itemize}
    \item \(x_B^G\): generic-oriented effort.
    \item \(x_B^O\): original-drug-oriented effort.
    \item \(C_B(\cdot)\): joint convex cost with resource rivalry across directions.
\end{itemize}

\vspace{0.15cm}
\textbf{Discipline condition:}
\begin{itemize}
    \item MAH does not directly shift \(\Psi_B^G\).
    \item MAH shifts \(\Psi_B^O\) via external implementation-route value.
\end{itemize}
\end{frame}

%------------------------------------------------
\begin{frame}{10. Original-Drug Authorization-Production Assignment}
\small
For original-drug projects generated under the type-B organizational state, the firm chooses an implementation route:
\[
H_O=\{ext,\; int,\; tr,\; ab\}.
\]

External route value:
\[
H_{B,O}^{ext}(z,p_m;M)=\lambda_{B,O}^{ext}(z,p_m)R_{B,O}(z)-p_m-\tau_{ext}(M).
\]

Internal route value:
\[
H_{B,O}^{int}(z)=\lambda_{int}(z)R_{B,O}(z)-\kappa_I(z).
\]

Per-project original-drug value:
\[
\Psi_B^O(z;p_m,M)=
\max\left\{
H_{B,O}^{ext},\,
H_{B,O}^{int},\,
H_{B,O}^{tr},\,
0
\right\}.
\]

\begin{itemize}
    \item MAH enters directly only through \(H_{B,O}^{ext}\).
    \item The internal route means the project is implemented under an integrated mode; it need not mean the whole firm is literally reclassified for all products.
\end{itemize}
\end{frame}

%------------------------------------------------
\begin{frame}{11. Why MAH Raises Original-Drug Effort in the Type-B State}
\small
Define the original-drug effort share under the type-B state:
\[
\omega_B^O(z;M,p_m)=
\frac{x_B^{O*}}{x_B^{G*}+x_B^{O*}}.
\]

Define the external implementation value of an original-drug idea:
\[
H_{B,O}^{ext}(z,p_m;M)=
\lambda_{B,O}^{ext}(z,p_m)R_{B,O}(z)-p_m-\tau_{ext}(M).
\]

\begin{block}{Main logic}
\[
\tau_{ext}\downarrow
\;\Rightarrow\;
H_{B,O}^{ext}\uparrow
\;\Rightarrow\;
\Psi_B^O\uparrow
\;\Rightarrow\;
x_B^{O*}\uparrow,\;\omega_B^O\uparrow
\]
\end{block}

\begin{itemize}
    \item The value of implementing original-drug ideas through the external route rises first.
    \item Type-B firms then adjust both effort level \((x_B^{O*})\) and effort direction \((\omega_B^O)\).
    \item MAH raises type-B original-drug effort through \textbf{external implementation value}, not direct scientific productivity.
\end{itemize}
\end{frame}

%------------------------------------------------
\begin{frame}{12. External Implementation Value: Object Definition}
\small
\textbf{Model payoff object:}
\[
H_{B,O}^{ext}(z,p_m;M)=
\lambda_{B,O}^{ext}(z,p_m)R_{B,O}(z)-p_m-\tau_{ext}(M).
\]

\textbf{Where:}
\begin{itemize}
    \item \(\lambda_{B,O}^{ext}(z,p_m)\in[0,1]\): reduced-form probability that a type-B original-drug idea is realized through the qualified external implementation route;
    \item \(R_{B,O}(z)\): private payoff after successful original-drug realization;
    \item \(p_m\): price of qualified external production service;
    \item \(\tau_{ext}(M)\): governance and compliance friction of the external implementation route;
    \item \(H_{B,O}^{ext}\): expected net payoff for a type-B state firm from realizing an original-drug idea through the external route.
\end{itemize}

\vspace{0.1cm}
\textbf{Interpretation:} this is a structural payoff object in the model, not an empirical reduced-form indicator.
\end{frame}

%------------------------------------------------
\begin{frame}{13. Comparative Statics: What Moves with MAH?}
\small
\begin{block}{Direct effect at fixed \((z,p_m)\)}
\[
\tau_{ext}\downarrow
\Rightarrow
H_{B,O}^{ext}\uparrow
\Rightarrow
\Psi_B^O\uparrow.
\]
\end{block}

\begin{itemize}
    \item \textbf{Route margin:} the external-route choice set expands.
    \item \textbf{Direction margin:} \(x_B^{O*}\uparrow\); and \(\omega_B^O\uparrow\) if \(\Psi_B^O\) rises relative to \(\Psi_B^G\).
    \item \textbf{Entry margin:} \(V_B\uparrow\), so research-oriented entry becomes more attractive.
    \item \textbf{GE caveat:} producer-side type-C profits are sign-ambiguous because \(p_m\), entry, and costs adjust endogenously.
\end{itemize}

\vspace{0.1cm}
\textbf{Maintained condition:} MAH primarily raises original-drug external implementation value, while generic payoff is unchanged or rises less.
\end{frame}

%------------------------------------------------
\begin{frame}{14. Dynamic Structure and Equilibrium}
\small
Firms can stay, switch dominant organizational state, or exit.

Bellman equation:
\[
V_s(z;M)=
\max\left\{
0,\;
\max_{j\in\mathcal{J}(s)}
W_{sj}(z;M)
\right\},
\qquad s\in\{A,B,C\}.
\]

where
\[
W_{sj}(z;M)=
\Pi_j(z;M,w,p_m)
-\kappa_{sj}(z;M)
+\beta E[V_j(z';M)\mid z,j].
\]

Contract-manufacturing market clearing:
\[
S_m(p_m,w;M)=D_m(p_m;M).
\]

\begin{itemize}
    \item Entry, exit, switching, and prices are all endogenous.
    \item We track both mechanism objects (route assignment, direction choice) and final innovation outcomes.
\end{itemize}
\end{frame}

%------------------------------------------------
\begin{frame}{15. Main Predictions}
\begin{block}{First-layer (main outcomes)}
\begin{itemize}
    \item external implementation-route use for original-drug projects rises,
    \item original-drug-oriented effort \(x_B^{O*}\) and share \(\omega_B^O\) rise,
    \item realized original-drug outcomes rise: \(N_O\), \(S_O\), and conversion \(Conv_O\).
\end{itemize}
\end{block}

\begin{block}{Second-layer (supporting margins)}
\begin{itemize}
    \item research-oriented entry margin becomes more attractive,
    \item organizational-state transitions \((A,B,C)\) may adjust,
    \item producer-side type-C effects are not guaranteed monotone in general equilibrium.
\end{itemize}
\end{block}
\vspace{0.1cm}
\textbf{Interpretation:} We prioritize innovation outcomes first, then use reallocation margins as mechanism validation.
\end{frame}

%------------------------------------------------
\begin{frame}{16. Empirical Mapping and How We Proceed}
\begin{itemize}
    \item \textbf{Module A (current approval-side data):}
    \begin{itemize}
        \item holder-producer split as route-use proxy,
        \item realized original-drug counts and shares,
        \item baseline reduced-form and baseline moments.
    \end{itemize}
    \item \textbf{Module B/C (future upgrades):}
    \begin{itemize}
        \item application-to-approval conversion and explicit entrusted production,
        \item capability-state construction, firm entry, and richer dominant-mode mapping.
    \end{itemize}
    \item Empirical claims remain modular: stronger data tighten the same model objects.
    \item Type mapping is approximate because \(A/B/C\) are organizational states, not literal legal firm categories.
\end{itemize}

\vspace{0.15cm}
\textbf{How we proceed (for this project stage):}
\begin{itemize}
    \item finish Module-A baseline facts and reduced forms,
    \item solve stationary equilibrium and estimate baseline moments,
    \item add Module-B/C data for conversion, capability, and entry,
    \item run counterfactuals after full moment expansion.
\end{itemize}
\end{frame}

%------------------------------------------------
\begin{frame}{17. Data-Estimation GAP: What We Have vs. What We Need}
\small
\begin{block}{What current Module-A data identifies well}
\begin{itemize}
    \item holder-producer split (route-use proxy),
    \item realized original-drug counts and shares,
    \item baseline pre/post moments for reduced-form discipline.
\end{itemize}
\end{block}

\begin{block}{Main gap for structural parameter estimation}
\begin{itemize}
    \item \textbf{Conversion denominator missing:} approval-only data cannot pin down application-to-approval conversion moments.
    \item \textbf{Capability state missing:} limited pre-reform firm capability measures for identifying state transition and dominant-mode mapping.
    \item \textbf{Route observability gap:} split proxy exists, but explicit entrusted-production indicators are still limited.
\end{itemize}
\end{block}

\textbf{Near-term fix:} keep baseline estimation modular now, then expand the moment set when Module-B/C merges are complete.
\end{frame}

%------------------------------------------------
\begin{frame}{18. What We Need Feedback On}
\begin{itemize}
    \item Is the object hierarchy clear enough: assignment as mechanism, original-drug innovation as final object?
    \item Is the direction-choice block \((x_B^G,x_B^O)\) the right minimal extension at current complexity?
    \item Is the organizational-state interpretation of \(A/B/C\) disciplined and empirically workable enough for the baseline?
    \item Is the Module A to Module B/C empirical roadmap convincing for identification?
\end{itemize}

\vspace{0.35cm}
\begin{block}{Bottom line}
MAH matters because it changes \emph{how} ideas become products and \emph{which} ideas firms choose to pursue.
\end{block}
\end{frame}

\end{document}
